% Generated draft: requires mathematical and source review.
\documentclass[11pt]{article}
\usepackage[margin=25mm]{geometry}
\usepackage{fontspec}
\setmainfont{Latin Modern Roman}
\usepackage{amsmath,amssymb,mathtools,unicode-math}
\setmathfont{Latin Modern Math}
\usepackage{enumitem,needspace,xurl,hyperref}
\setlength{\parindent}{0pt}
\setlength{\parskip}{5pt}
\setlength{\emergencystretch}{3em}
\newcommand{\N}{\mathbb{N}}
\newcommand{\Z}{\mathbb{Z}}
\newcommand{\Q}{\mathbb{Q}}
\newcommand{\R}{\mathbb{R}}
\newcommand{\C}{\mathbb{C}}
\newcommand{\F}{\mathbb{F}}
\newcommand{\A}{\mathbb{A}}
\newcommand{\PP}{\mathbb{P}}
\newcommand{\E}{\mathbb{E}}
\newcommand{\eps}{\varepsilon}
\newcommand{\dd}{\,\mathrm{d}}
\begin{document}
\section*{UnsolvedMath ID: 700781}
\textbf{Source problem code:} catalogue entry
\subsection{Definitions and mathematical statement}
Let \(k\) be an infinite cardinal, and let \(I\) be a set with \(|I|=k\). Define the Cartesian power of the additive group of integers by
\[
P_k=\prod_{i\in I}\Z,
\]
with coordinatewise addition; its elements are all functions \(I\to\Z\), without a finite-support restriction. An epimorphic image means an abstract group \(G\) admitting a surjective group homomorphism \(P_k\twoheadrightarrow G\). Every such \(G\) is abelian.

For each infinite cardinal \(k\), describe, up to group isomorphism, all epimorphic images of \(P_k\). Equivalently, characterize the abelian groups \(G\) for which there exists a subgroup \(H\leq P_k\) such that
\[
G\cong P_k/H.
\]
The task is to give a description of these groups beyond this defining quotient formulation. No topology or continuity condition is specified in the source.

The supplied source states that such a description is known for \(k=\omega_0=\aleph_0\), citing R. Nunke (1963). This is retained as a source-reported qualification and has not been independently verified.
\subsection{Short English statement}
For each infinite cardinal \(k\), classify the groups obtainable as surjective homomorphic images of the full Cartesian product of \(k\) copies of \(\Z\). The source reports that the countably infinite case was described by R. Nunke.
\subsection{Sources}
\begin{itemize}
\item Local UnsolvedMath registry, supplied record with permanent ID \(700781\), exact source code \texttt{catalogue entry}, title \emph{Kourovka Notebook Problem 11.89}. The registry labels the problem open; this label is supplied metadata, not an independently verified literature assessment.
\item \emph{The Kourovka Notebook (Unsolved Problems in Group Theory)}, Problem 11.89, 11th issue (1990); proposed by S. V. Rychkov. The supplied background also identifies the 21st issue edition (2026), \texttt{arXiv:1401.0300v48}. The registry collection names E. I. Khukhro and V. D. Mazurov as editors.
\item R. Nunke, \emph{Acta Sci. Math. (Szeged)} \textbf{23} (1963), 67--73. Supplied as the reference for the description when \(k=\omega_0\); no article title is provided.
\end{itemize}
\end{document}
